\section{Conclusion and future work}
\label{sec:conclusion}

This thesis aims to study the generation and comprehension of referring expressions in language games.
Hereby, two agents, a sender and a receiver, are tasked to solve tasks together by learning and communicating messages in an emerged language.
More specifically, the agents learn to ground referring expressions that describe objects in visual scenes.
As the agents will exploit all present biases in the datasets, three new artificial datasets are created where these biases are precisely controlled for.
Objects in the visual scene have three attributes, \emph{shape}, \emph{size} and \emph{color}.
The target object is uniquely identifiable, while distractors have at least one differing attribute.
While the 'Dale-2' dataset contains only one distractor, the 'Dale-5' dataset contains four distractors.
In the 'CLEVR color' dataset, the \emph{color} is the only attribute that is different, while all other attributes are shared between the target object and the distractors.
The agents are expected to use these specific patterns to refer to a target object and respectively comprehend the generated referring expression.
All unwanted bias is reduced by randomizing the creation of the visual scenes, so that the neural models can only use the attributes instead.

To validate that neural models are able to use this controlled bias, several neural models are trained to solve tasks on these datasets.
The results on the \emph{object identification} task show that the models are able to comprehend referring expressions and ground them in bounding boxes, as they can consistently point to the correct target object.
In the \emph{referring expression generation} tasks, the models can successfully extract visual features and generate referring expressions to discriminate the target object from the distractors.
While the task is solved very well when presented with bounding boxes, the models struggle more when presented with the complete scene at once, especially when many objects are present.
However, the scores are far above random guesses and show that the introduced bias in the datasets can be exploited.
Finally, the models learn to combine geometric information with the extracted visual features in the \emph{reference resolution} task.
The models only manage to predict precise coordinates of the target object given a referring expression when two objects are present.
On the other side, when tasked to attend to rougher regions around the target object instead, the models are able to solve the task almost perfectly on all datasets.
This shows again that the bias in the dataset is enough to ground referring expressions in a spatial context.

In the next step, the two agents are trained on the same tasks.
The increased complexity in the architecture of the language games shows in the results.
The agents have more difficulties to solve the tasks.
However, in almost all tasks, the agents manage to ground some information in referring expressions of the emerged language.
The datasets and the tasks have hereby the highest influence on the success of grounding referring expressions.
The simpler 'Dale-2' dataset allows the agents to ground referring expressions the best in their messages.
While the agents are still able to solve the tasks in many language games on the 'Dale-5' dataset, they have the biggest problems on the 'CLEVR color' dataset.
The main reason seems to be the increased number of distractors in combination with extracting features from the whole visual scene as opposed to bounding boxes for each object.
Additionally, the results suggest a difference of difficulty to ground each attributes.
The \emph{shape} and \emph{size} seem to be grounded easier than the \emph{color}.

Finally, looking closer at the emerged languages, it can be seen that the used referring expressions don't align perfectly with any referring expression that would be common in English.
Interestingly, as common in English referring expressions, the emerged languages seem to assign a salience order to the attributes.
On the 'Dale' dataset, the \emph{size} seems to be the most salient attribute, followed by the \emph{shape} and lastly the \emph{color}.
While the order is different from one used in English, the alignment can be influenced by introducing English referring expressions during training.
Then, the \emph{shape} becomes the most salient attribute.
For the 'CLEVR color' dataset, unsurprisingly, the most salient attribute is the \emph{color}.
However, the agents seem to utilize some other patterns apart from the three attributes.

This thesis lies the basis for several future experiments.
As discussed, some attributes seem to be grounded easier than others.
Future work should focus on the reasons for this.
For instance, the agents can be trained on new datasets like the 'CLEVR color' dataset where the only differing attributes are \emph{shape} and \emph{size} respectively.

Secondly, the emerged languages can be studied in more depth.
While this thesis focuses on an alignment with English referring expression, other linguistic properties like composition could give better insights in how the referring expressions are structured.

Furthermore, at the moment, the sender has the same architecture across all experiments, and the adpated receiver is mainly responsible for the differences in the referring expressions.
It would be interesting however, how the referring expressions change if the receiver is kept stable, but the sender is variable.
This could give insights about the influence that each agent has on the structure of referring expressions.

Finally, communication only works in one direction, from the sender to the receiver.
The emergence of referring expressions is therefore solely based on the final training loss of the receiver.
It might support the agents to have a bidirectional communication in which the receiver can respond, and the agents can agree directly on referring expressions.